\chapter{Control Flow}
\label{ch:control-flow}

Control flow constructs determine the order in which statements are executed. Novus provides a comprehensive set of control flow mechanisms, from basic conditionals to sophisticated pattern matching. This chapter covers all the ways you can direct program execution in Novus.

\section{Conditional Statements}
\label{sec:conditionals}

\subsection{if and else}
\label{subsec:if-else}

The \texttt{if} statement executes code based on a boolean condition:

\begin{lstlisting}
fn classify_temperature(temp: i32) -> i32 {
    if temp < 0 {
        return -1  // Freezing
    } else if temp < 20 {
        return 0   // Cold
    } else if temp < 30 {
        return 1   // Comfortable
    } else {
        return 2   // Hot
    }
}
\end{lstlisting}

The condition must be of type \texttt{bool}. Unlike C, Novus does not implicitly convert integers to booleans---you must use explicit comparisons.

\textbf{Blocks are required.} Unlike C, the braces around the body are mandatory, even for single statements:

\begin{lstlisting}
// Correct
if x > 0 {
    return x
}

// Syntax error: missing braces
// if x > 0 return x
\end{lstlisting}

\subsection{if as an Expression}
\label{subsec:if-expression}

In Novus, \texttt{if} can be used as an expression when both branches return a value:

\begin{lstlisting}
let max = if a > b { a } else { b }

// Also valid with blocks
let result = if condition {
    compute_something()
} else {
    compute_fallback()
}
\end{lstlisting}

When used as an expression, both the \texttt{if} and \texttt{else} branches must be present and must return the same type:

\begin{lstlisting}
// Error: branches return different types
let bad = if condition { 42 } else { "hello" }
\end{lstlisting}

\subsection{Short-Circuit Evaluation}
\label{subsec:short-circuit}

The logical operators \texttt{\&\&} (and) and \texttt{||} (or) use short-circuit evaluation:

\begin{itemize}
    \item \texttt{a \&\& b}: If \texttt{a} is false, \texttt{b} is not evaluated
    \item \texttt{a || b}: If \texttt{a} is true, \texttt{b} is not evaluated
\end{itemize}

This allows safe idioms like:

\begin{lstlisting}
// ptr is only dereferenced if non-null
if ptr != 0 && (*ptr).value > 0 {
    // Safe to use ptr here
}

// check_signal is only called if running is true
if running && check_signal() {
    handle_signal()
}
\end{lstlisting}

\subsection{if let}
\label{subsec:if-let}

The \texttt{if let} construct checks whether a value is non-null (for pointers) or non-zero (for integers) and binds it to a variable. This is primarily used when working with raw pointers from AmigaOS FFI calls:

\begin{lstlisting}
fn process_pointer(ptr: *u8) -> i32 {
    if let p = ptr {
        // p is bound to ptr and is guaranteed non-null here
        return (i32)p
    } else {
        // ptr was null
        return -1
    }
}
\end{lstlisting}

\textbf{Note:} For pattern matching on \texttt{Option<T>} and \texttt{Result<T, E>} types, use \texttt{let-else} (Section~\ref{subsec:let-else}) or \texttt{match} (Section~\ref{sec:match}) instead.

The \texttt{if let} construct is particularly useful with pointers returned from AmigaOS functions that may return null on failure:

\begin{lstlisting}
fn open_library(name: *u8) -> i32 {
    let lib = OpenLibrary(name, 0)

    if let base = lib {
        // Library opened successfully, base is non-null
        CloseLibrary(base)
        return 0
    } else {
        // Library not found
        return 1
    }
}
\end{lstlisting}

\subsection{let-else}
\label{subsec:let-else}

The \texttt{let-else} statement handles pattern matching where the \texttt{else} block must diverge (i.e., return, break, continue, or panic). This is ideal for early-exit patterns:

\begin{lstlisting}
enum MaybeValue {
    Some(i32),
    None
}

fn process(input: MaybeValue) -> i32 {
    let MaybeValue::Some(value) = input else {
        // This block runs if pattern doesn't match
        // Must diverge - cannot fall through
        return -1
    }

    // value is now bound and available here
    return value * 2
}
\end{lstlisting}

The \texttt{else} block \emph{must} diverge. The compiler will reject code where the else block could fall through to the subsequent statements.

This pattern works well with \texttt{Result} and \texttt{Option} types:

\begin{lstlisting}
from std::core import Result
from std::error::errors import DosError

fn read_file(path: *u8) -> Result<i32, DosError> {
    let Result::Ok(handle) = open_file(path) else {
        return Result::Err(DosError::NotFound)
    }

    // handle is now available for use
    let bytes = read_data(handle)
    close_file(handle)
    return Result::Ok(bytes)
}
\end{lstlisting}

\section{Loops}
\label{sec:loops}

Novus provides several loop constructs, each suited to different use cases.

\subsection{while Loops}
\label{subsec:while}

The \texttt{while} loop executes its body as long as the condition is true:

\begin{lstlisting}
var count = 0
while count < 10 {
    count++
}
// count is now 10
\end{lstlisting}

The condition is evaluated before each iteration. If the condition is initially false, the body never executes.

\begin{lstlisting}
fn find_first_even(start: i32, max: i32) -> i32 {
    var n = start
    while n < max {
        if n % 2 == 0 {
            return n  // Found an even number
        }
        n++
    }
    return -1  // None found
}
\end{lstlisting}

\subsection{forever Loops}
\label{subsec:forever}

The \texttt{forever} loop creates an infinite loop that must be exited with \texttt{break} or \texttt{return}:

\begin{lstlisting}
fn wait_for_signal() -> i32 {
    var iterations = 0
    forever {
        iterations++
        if check_signal() {
            break  // Exit the loop
        }
        if iterations > 1000 {
            return -1  // Timeout
        }
    }
    return iterations
}
\end{lstlisting}

The \texttt{forever} keyword is clearer than \texttt{while true} and signals the programmer's intent that the loop is designed to be exited via control flow rather than a condition becoming false.

This is particularly useful for event loops in Amiga applications:

\begin{lstlisting}
fn main_event_loop(window: *Window) {
    forever {
        let msg = GetMsg(window.UserPort)
        if let m = msg {
            let class = m.Class
            ReplyMsg(m)

            if class == IDCMP_CLOSEWINDOW {
                break  // User closed the window
            }

            handle_message(class)
        }

        WaitPort(window.UserPort)
    }
}
\end{lstlisting}

\subsection{C-Style for Loops}
\label{subsec:for-cstyle}

Novus supports traditional C-style \texttt{for} loops with initialization, condition, and update expressions:

\begin{lstlisting}
// Sum numbers 0 through 9
var sum = 0
for var i = 0; i < 10; i++ {
    sum = sum + i
}
// sum is now 45
\end{lstlisting}

The initialization can declare a new variable (with \texttt{var} or \texttt{let}) or assign to an existing one. The loop variable declared in the initialization is scoped to the loop.

\begin{lstlisting}
// Skip even numbers with continue
var odd_sum = 0
for var j = 0; j < 10; j++ {
    if j % 2 == 0 {
        continue  // Skip even numbers
    }
    odd_sum = odd_sum + j
}
// odd_sum is 1 + 3 + 5 + 7 + 9 = 25
\end{lstlisting}

\subsection{Range-Based for Loops}
\label{subsec:for-range}

Novus provides syntactic sugar for iterating over numeric ranges:

\begin{lstlisting}
// Exclusive range: 0, 1, 2, 3, 4 (excludes 5)
var sum = 0
for i in 0..5 {
    sum = sum + i
}
// sum is 10
\end{lstlisting}

\begin{lstlisting}
// Inclusive range: 1, 2, 3, 4, 5 (includes 5)
var sum = 0
for i in 1..=5 {
    sum = sum + i
}
// sum is 15
\end{lstlisting}

Ranges can use variables:

\begin{lstlisting}
let start = 2
let end = 7
var sum = 0
for i in start..end {
    sum = sum + i
}
// sum is 2 + 3 + 4 + 5 + 6 = 20
\end{lstlisting}

An empty range (where start equals end) executes zero iterations:

\begin{lstlisting}
var count = 0
for i in 5..5 {
    count++  // Never executes
}
// count is 0
\end{lstlisting}

\subsection{for-in Loops with Collections}
\label{subsec:for-in}

The \texttt{for-in} loop iterates over collections that provide \texttt{get()} and \texttt{len()} methods:

\begin{lstlisting}
// ... (with Vec imported from std::collections::vec)
let vec: Vec<i32> = Vec { [10, 20, 30, 40, 50] }

var sum: i32 = 0
for value in vec {
    sum += value
}
// sum is 150
\end{lstlisting}

The loop variable receives a copy of each element. The collection itself remains usable after the loop completes. For collections containing non-copyable types, the iteration behavior depends on the collection's implementation.

Any type that provides \texttt{get(index: i32)} and \texttt{len()} methods can be used with \texttt{for-in}. The compiler generates code equivalent to a C-style for loop that calls these methods.

You can declare the loop variable as mutable if you need to modify the local copy:

\begin{lstlisting}
for var item in collection {
    item = transform(item)  // Modifies local copy
    process(item)
}
\end{lstlisting}

\section{Loop Control Statements}
\label{sec:loop-control}

\subsection{break and continue}
\label{subsec:break-continue}

The \texttt{break} statement immediately exits the innermost loop:

\begin{lstlisting}
var i = 0
while i < 100 {
    if i == 42 {
        break  // Exit immediately
    }
    i++
}
// i is 42
\end{lstlisting}

The \texttt{continue} statement skips the rest of the current iteration and proceeds to the next:

\begin{lstlisting}
var sum = 0
var i = 0
while i < 10 {
    i++
    if i == 5 {
        continue  // Skip adding 5
    }
    sum = sum + i
}
// sum is 1+2+3+4+6+7+8+9+10 = 50 (5 is skipped)
\end{lstlisting}

\subsection{Labeled Loops}
\label{subsec:labeled-loops}

When working with nested loops, you can label loops and use those labels with \texttt{break} and \texttt{continue} to control outer loops:

\begin{lstlisting}
var count = 0
outer: while true {
    var inner = 0
    while inner < 5 {
        count++
        if count == 7 {
            break outer  // Exit the outer loop
        }
        inner++
    }
}
// count is 7
\end{lstlisting}

Labels work with all loop types:

\begin{lstlisting}
outer: forever {
    for i in 0..10 {
        for j in 0..10 {
            if i * j > 50 {
                break outer  // Exit all three loops
            }
        }
    }
}
\end{lstlisting}

The \texttt{continue} statement also works with labels:

\begin{lstlisting}
var sum = 0
var i = 0
outer: while i < 5 {
    i++
    var j = 0
    while j < 3 {
        j++
        if i == 3 && j == 2 {
            continue outer  // Skip to next i iteration
        }
        sum++
    }
}
\end{lstlisting}

\textbf{Note:} Unlike C labels, Novus labels are only valid for loop control. Novus does not support \texttt{goto} statements---labels can only be used with \texttt{break} and \texttt{continue}.

\subsection{Postfix Conditions}
\label{subsec:postfix-conditions}

Novus supports postfix \texttt{if} and \texttt{unless} for concise conditional statements:

\begin{lstlisting}
fn classify(x: i32) -> i32 {
    return -1 if x < 0    // Return -1 if negative
    return 1 if x > 0     // Return 1 if positive
    return 0              // Zero
}
\end{lstlisting}

The \texttt{unless} keyword is the logical inverse of \texttt{if}:

\begin{lstlisting}
fn check(valid: bool) -> i32 {
    return -1 unless valid  // Return -1 if NOT valid
    return 0
}
\end{lstlisting}

Postfix conditions work with \texttt{break} and \texttt{continue}:

\begin{lstlisting}
var sum = 0
var m = 0
while m < 10 {
    m++
    continue if m % 3 == 0  // Skip multiples of 3
    sum = sum + m
}
// sum is 1+2+4+5+7+8+10 = 37

outer: while true {
    var k = 0
    while k < 100 {
        k++
        break outer if k == 5  // Exit outer when k reaches 5
    }
}
\end{lstlisting}

\section{Pattern Matching with match}
\label{sec:match}

The \texttt{match} expression provides exhaustive pattern matching over values. It's more powerful than a simple switch statement and is the idiomatic way to handle enums in Novus.

\subsection{Basic match Expressions}
\label{subsec:match-basics}

\begin{lstlisting}
fn http_status_category(code: i32) -> i32 {
    match code {
        200 => return 1,  // Success
        201 => return 1,
        204 => return 1,
        400 => return 2,  // Client Error
        401 => return 2,
        404 => return 2,
        500 => return 3,  // Server Error
        503 => return 3,
        _   => return 0   // Unknown (wildcard)
    }
}
\end{lstlisting}

The underscore (\texttt{\_}) is a wildcard pattern that matches any value not handled by previous arms.

\begin{comingfromc}
\texttt{match} is similar to C's \texttt{switch}, but safer and more powerful:

\begin{itemize}
\item \textbf{No fallthrough}: Each arm is independent; no \texttt{break} needed
\item \textbf{Exhaustive}: The compiler ensures all cases are handled
\item \textbf{Expressions}: \texttt{match} returns a value, not just executes code
\item \textbf{Patterns}: Can destructure enums, tuples, and structs
\end{itemize}

\begin{tabular}{ll}
\textbf{C switch} & \textbf{Novus match} \\
\texttt{case 1: ...; break;} & \texttt{1 => ...,} \\
\texttt{default: ...} & \texttt{\_ => ...} \\
\end{tabular}

Forget \texttt{break}---the absence of fallthrough eliminates a whole class of bugs.
\end{comingfromc}

\subsection{Matching on Enums}
\label{subsec:match-enums}

Pattern matching is essential for working with enums, especially those with associated data:

\begin{lstlisting}
enum Status {
    Success,
    Pending,
    Error(i32)
}

fn handle_status(status: Status) -> i32 {
    match status {
        Status::Success => return 0,
        Status::Pending => return 1,
        Status::Error(code) => return code  // Bind the error code
    }
}
\end{lstlisting}

When an enum variant carries data, the pattern can bind that data to a variable for use in the match arm.

\subsection{Match with Generic Enums}
\label{subsec:match-generic}

Pattern matching works seamlessly with generic enums like \texttt{Option<T>} and \texttt{Result<T, E>}:

\begin{lstlisting}
from std::core import Option

fn get_value_or_default(opt: Option<i32>) -> i32 {
    match opt {
        Option::Some(val) => return val,
        Option::None => return 0
    }
}
\end{lstlisting}

\begin{lstlisting}
from std::core import Result
from std::error::errors import DosError

fn process_result(res: Result<i32, DosError>) -> i32 {
    match res {
        Result::Ok(value) => return value,
        Result::Err(_) => return -1  // Ignore error details
    }
}
\end{lstlisting}

\subsection{Match with Blocks}
\label{subsec:match-blocks}

When a match arm requires multiple statements, use a block:

\begin{lstlisting}
fn sign(n: i32) -> i32 {
    match n {
        0 => return 0,
        _ => {
            if n > 0 {
                return 1   // Positive
            }
            return -1      // Negative
        }
    }
}
\end{lstlisting}

\subsection{Match with Guards}
\label{subsec:match-guards}

\begin{tcolorbox}[colback=red!5!white,colframe=red!75!black,title={\textbf{Work in Progress}}]
Match guards are parsed but code generation has known issues. Use \texttt{if}/\texttt{else} chains as an alternative until this is resolved.
\end{tcolorbox}

Match arms can include guard clauses with \texttt{if} to add additional conditions:

\begin{lstlisting}
fn classify_number(n: i32) -> i32 {
    match n {
        x if x < 0 => return -1,      // Negative
        0 => return 0,                 // Zero
        x if x <= 10 => return 1,     // Small positive
        _ => return 2                  // Large positive
    }
}
\end{lstlisting}

The guard is evaluated after the pattern matches but before the arm is executed. If the guard evaluates to false, matching continues with the next arm.

\subsection{Matching Integer Literals}
\label{subsec:match-integers}

Novus supports matching on decimal, hexadecimal, and binary integer literals:

\begin{lstlisting}
fn parse_permission(bits: i32) -> i32 {
    match bits {
        %111 => return 7,  // rwx (binary)
        %110 => return 6,  // rw-
        %100 => return 4,  // r--
        _    => return 0
    }
}

fn hex_component(hex: i32) -> i32 {
    match hex {
        $00 => return 0,    // Hexadecimal
        $FF => return 255,
        $80 => return 128,
        _   => return hex
    }
}
\end{lstlisting}

\subsection{Exhaustiveness Checking}
\label{subsec:match-exhaustiveness}

The Novus compiler verifies that match expressions are exhaustive---every possible value must be handled. For enums, this means every variant must have a corresponding arm (or a wildcard):

\begin{lstlisting}
enum Direction {
    North,
    South,
    East,
    West
}

fn direction_to_int(dir: Direction) -> i32 {
    match dir {
        Direction::North => return 0,
        Direction::South => return 1,
        Direction::East => return 2
        // Error: non-exhaustive pattern, Direction::West not covered
    }
}
\end{lstlisting}

Either add a case for \texttt{Direction::West} or use a wildcard pattern.

\section{Error Propagation with the Try Operator}
\label{sec:try-operator}

The try operator (\texttt{?}) provides concise error propagation for functions returning \texttt{Result<T, E>}.

\subsection{Basic Usage}
\label{subsec:try-basic}

Without the try operator, error handling requires verbose match expressions:

\begin{lstlisting}
fn process_file_manual(path: *u8) -> Result<i32, DosError> {
    let handle_result = open_file(path)
    let handle = match handle_result {
        Result::Ok(h) => h,
        Result::Err(e) => return Result::Err(e)
    }

    let bytes_result = read_file(handle)
    let bytes = match bytes_result {
        Result::Ok(b) => b,
        Result::Err(e) => return Result::Err(e)
    }

    close_file(handle)
    return Result::Ok(bytes)
}
\end{lstlisting}

With the try operator, this becomes:

\begin{lstlisting}
fn process_file(path: *u8) -> Result<i32, DosError> {
    let handle = open_file(path)?
    let bytes = read_file(handle)?
    close_file(handle)?

    return Result::Ok(bytes)
}
\end{lstlisting}

The \texttt{?} operator:
\begin{enumerate}
    \item Unwraps \texttt{Result::Ok(value)} and yields \texttt{value}
    \item On \texttt{Result::Err(e)}, immediately returns \texttt{Result::Err(e)} from the current function
\end{enumerate}

\subsection{Try Operator in Expressions}
\label{subsec:try-expressions}

The try operator can be used within expressions:

\begin{lstlisting}
fn get_file_size(path: *u8) -> Result<i32, DosError> {
    let handle = open_file(path)?
    let bytes = read_file(handle)?
    return Result::Ok(bytes + 100)  // Add header overhead
}
\end{lstlisting}

\subsection{Chaining Operations}
\label{subsec:try-chaining}

The try operator excels at chaining multiple fallible operations:

\begin{lstlisting}
fn complex_operation() -> Result<i32, DosError> {
    let file = open_file("data.txt")?
    let header = read_header(file)?
    let body = read_body(file, header.length)?
    let result = process_data(body)?
    close_file(file)?

    return Result::Ok(result)
}
\end{lstlisting}

Each \texttt{?} acts as an early-exit point---if any operation fails, the error propagates immediately to the caller.

\subsection{Combining with defer}
\label{subsec:try-defer}

The try operator pairs well with \texttt{defer} for resource cleanup:

\begin{lstlisting}
fn safe_file_operation(path: *u8) -> Result<i32, DosError> {
    let handle = open_file(path)?
    defer {
        close_file(handle)
    }

    let header = read_header(handle)?   // If this fails...
    let body = read_body(handle)?       // ...or this...

    // The defer block still runs, closing the file
    return Result::Ok(process(body))
}
\end{lstlisting}

\section{The defer Statement}
\label{sec:defer}

The \texttt{defer} statement schedules code to run when the current scope exits, regardless of how it exits (normal completion, early return, or error propagation).

\subsection{Basic Usage}
\label{subsec:defer-basic}

\begin{lstlisting}
fn example() -> i32 {
    let buffer = allocate_buffer(1024)
    defer {
        free_buffer(buffer)
    }

    // Use buffer...
    if some_condition() {
        return -1  // defer block still runs
    }

    return 0  // defer block runs before returning
}
\end{lstlisting}

\subsection{Multiple Defers}
\label{subsec:defer-multiple}

Multiple \texttt{defer} statements execute in reverse order (last-in, first-out):

\begin{lstlisting}
fn multi_resource() {
    let resource1 = acquire_first()
    defer { release_first(resource1) }

    let resource2 = acquire_second()
    defer { release_second(resource2) }

    let resource3 = acquire_third()
    defer { release_third(resource3) }

    // Use resources...

    // On exit: release_third, then release_second, then release_first
}
\end{lstlisting}

This LIFO order ensures resources are released in the reverse order of acquisition, which is typically the correct order for dependent resources.

\subsection{Defer with AmigaOS Resources}
\label{subsec:defer-amiga}

The \texttt{defer} statement is particularly useful for managing AmigaOS resources:

\begin{lstlisting}
fn open_window_safe() -> i32 {
    let screen = LockPubScreen(0)
    if screen == 0 {
        return -1
    }
    defer { UnlockPubScreen(0, screen) }

    let visual_info = GetVisualInfo(screen, 0)
    if visual_info == 0 {
        return -2
    }
    defer { FreeVisualInfo(visual_info) }

    let window = OpenWindowTags(screen)
    if window == 0 {
        return -3
    }
    defer { CloseWindow(window) }

    // Use window...
    process_events(window)

    return 0
    // Resources released in order: window, visual_info, screen
}
\end{lstlisting}

\section{The using Statement}
\label{sec:using}

The \texttt{using} statement provides scoped resource management for types that implement the \texttt{Drop} trait. When the block exits, the resource's \texttt{drop} method is called automatically.

\begin{lstlisting}
pub struct Resource {
    value: i32
}

impl Drop for Resource {
    fn drop(&var self) {
        // Cleanup happens here automatically
    }
}

fn create_resource(v: i32) -> Resource {
    return Resource { value: v }
}

fn process_data() -> i32 {
    using create_resource(42) {
        // The resource exists for this block
        let x: i32 = 10
        // ... use x ...
    }
    // Resource's drop() called automatically here

    return 0
}
\end{lstlisting}

The \texttt{using} statement takes an expression that returns a type implementing \texttt{Drop}. The resource is created, the block executes, and then \texttt{drop()} is called---regardless of how the block exits.

\subsection{using vs defer}
\label{subsec:using-vs-defer}

Both \texttt{using} and \texttt{defer} provide deterministic cleanup, but they serve different purposes:

\textbf{Use \texttt{defer}} when:
\begin{itemize}
    \item You need to bind the resource to a variable for use in the scope
    \item The resource must live for the entire function scope
    \item You need cleanup after early returns from anywhere in the function
    \item You're working with raw FFI resources that don't implement \texttt{Drop}
\end{itemize}

\textbf{Use \texttt{using}} when:
\begin{itemize}
    \item The resource is only needed to exist during a specific block
    \item You don't need to access the resource directly (e.g., a lock or guard)
    \item The resource type implements the \texttt{Drop} trait
\end{itemize}

Example showing the difference:

\begin{lstlisting}
fn process_files() -> i32 {
    // defer: resource is bound to a variable, cleanup at function exit
    let handle = open_file("log.txt")
    defer { close_file(handle) }

    // using: resource exists for block scope only
    // Useful for locks, guards, or temporary resources
    using acquire_lock() {
        // Lock is held for this block
        do_protected_work()
    }
    // Lock released here

    // We can still use handle
    write_to_file(handle, "Done")
    return 0
    // handle closed here by defer
}
\end{lstlisting}

\section{Summary}
\label{sec:control-flow-summary}

Novus provides a complete set of control flow constructs:

\begin{itemize}
    \item \textbf{Conditionals}: \texttt{if}/\texttt{else}, \texttt{if let}, \texttt{let-else} for pattern matching with early exit
    \item \textbf{Short-circuit evaluation}: \texttt{\&\&} and \texttt{||} operators for efficient boolean logic
    \item \textbf{Loops}: \texttt{while}, \texttt{forever}, C-style \texttt{for}, range-based \texttt{for}, and \texttt{for-in}
    \item \textbf{Loop control}: \texttt{break}, \texttt{continue}, labeled loops for nested control
    \item \textbf{Postfix conditions}: \texttt{if} and \texttt{unless} for concise conditional statements
    \item \textbf{Pattern matching}: \texttt{match} with exhaustiveness checking, guards, and literal patterns
    \item \textbf{Error propagation}: The \texttt{?} operator for concise \texttt{Result} handling
    \item \textbf{Resource management}: \texttt{defer} and \texttt{using} for deterministic cleanup
\end{itemize}

These constructs work together to enable clear, safe, and expressive code. The emphasis on explicit control flow and exhaustive pattern matching helps prevent bugs while keeping the code readable.

In the next chapter, we'll explore functions and methods, including function definitions, parameter passing, generic functions, and impl blocks.
